\chapter{Fundamentação teórica}
\label{cap:fundamentacao}

A conexão de empreendimentos com múltiplas unidades consumidoras ao sistema de distribuição requer a definição das condições de atendimento, do dimensionamento da entrada de energia e da organização da medição. Essas decisões são condicionadas pela regulamentação do setor elétrico e pelas especificações técnicas da distribuidora responsável pela área de concessão. No Ceará, os requisitos anteriormente estabelecidos pela especificação CNC-OMBR-MAT-18-0126-EDCE foram substituídos pelos da CNC-NDBR-DBR-25-1580-EDBR, documento de abrangência nacional no âmbito da Enel. A avaliação dessa transição exige distinguir disposições regulatórias, critérios construtivos e requisitos regionais.

\section{Regulamentação do fornecimento de energia elétrica}
\label{sec:regulamentacao}

As condições gerais de prestação do serviço público de distribuição são regulamentadas pela Agência Nacional de Energia Elétrica (ANEEL). A Resolução Normativa ANEEL nº 1.000/2021 consolida disposições sobre conexão, atendimento e responsabilidades dos agentes envolvidos. Os critérios relativos ao grupo e à tensão de fornecimento devem, portanto, ser interpretados conjuntamente com as normas técnicas da distribuidora, sem afastar as disposições regulatórias aplicáveis. A Resolução Normativa ANEEL nº 1.059/2023 promoveu alterações em dispositivos da regulamentação original, conforme \citeonline{aneel1000} e \citeonline{aneel1059}.

O enquadramento das unidades consumidoras nos Grupos A e B não decorre exclusivamente da potência demandada pelo empreendimento. De acordo com o art. 23 da Resolução Normativa ANEEL nº 1.000/2021, a carga instalada, as características da rede de distribuição e as hipóteses regulamentares de atendimento constituem parâmetros para definição da tensão de conexão. A referência geral de 75~kW para determinadas conexões em tensão inferior a 2,3~kV não elimina as possibilidades excepcionais previstas na regulamentação. Por essa razão, limites atribuídos às unidades consumidoras individuais não devem ser confundidos com os limites de demanda da entrada coletiva.

A entrada de energia está sujeita, ainda, aos requisitos de segurança e desempenho das instalações elétricas. A ABNT NBR 5410 e outros documentos técnicos referenciados pelas especificações da Enel integram esse referencial, respeitadas as respectivas áreas de aplicação. A análise documental realizada pela distribuidora não transfere ao agente responsável pelo serviço público a responsabilidade técnica do profissional habilitado pelo projeto e pelas instalações particulares.

\section{Conexão coletiva e múltiplas unidades consumidoras}
\label{sec:conexaocoletiva}

A conexão coletiva caracteriza-se pelo atendimento a um conjunto de unidades consumidoras mediante infraestrutura de entrada e distribuição concebida para o empreendimento. Em edifícios residenciais, os circuitos destinados às unidades autônomas coexistem com circuitos associados aos serviços de uso comum. Essa configuração requer a compatibilização dos pontos de conexão, da distribuição interna, dos dispositivos de proteção e das medições individualizadas.

Na especificação CNC-OMBR-MAT-18-0126-EDCE, o conceito de prédio de múltiplas unidades consumidoras (PMUC) orienta requisitos associados à concentração das medições e à alimentação dos consumidores vinculados à mesma edificação. A especificação abrange unidades autônomas e instalações das áreas comuns, estabelecendo condições próprias para a entrada de serviço, a apresentação do projeto e o arranjo dos equipamentos de proteção e medição, conforme \citeonline{enel0126}.

A CNC-NDBR-DBR-25-1580-EDBR utiliza a denominação \emph{conexão coletiva} e abrange empreendimentos horizontais e verticais. O item 7.8.1 distingue dessa condição as edificações contíguas ou geminadas sem área de uso comum e alinhadas à via pública, às quais se aplicam critérios de conexão individual. A delimitação da modalidade de atendimento depende, desse modo, das características funcionais e elétricas do conjunto, não apenas de sua configuração arquitetônica, conforme \citeonline[item~7.8.1]{enel1580}.

\section{Carga instalada e demanda no dimensionamento da entrada}
\label{sec:demandas}

A carga instalada corresponde à soma das potências nominais dos equipamentos previstos ou existentes na instalação, de acordo com os critérios aplicáveis a cada categoria de consumidor. A demanda expressa a solicitação de potência em determinada condição de operação. Em empreendimentos com várias unidades consumidoras, a simultaneidade de utilização dos equipamentos constitui parâmetro essencial do dimensionamento da entrada geral, uma vez que a soma aritmética das cargas individuais não representa, necessariamente, a maior potência requerida ao sistema.

A especificação CNC-OMBR-MAT-18-0126-EDCE apresenta, em sua seção 7.11, critérios de estimativa de demanda associados à área útil das unidades residenciais e procedimentos destinados a outras categorias de carga. A aplicação desses critérios depende da configuração das unidades e das características das cargas previstas. Na especificação CNC-NDBR-DBR-25-1580-EDBR, a determinação da demanda continua vinculada ao profissional habilitado, com referência às metodologias recomendadas pela distribuidora e aos documentos técnicos associados. Os dispositivos correspondentes estão apresentados na seção~7.11 da especificação 0126 e no item~7.3.2 da especificação 1580, conforme \citeonline{enel0126} e \citeonline{enel1580}.

A distinção entre carga instalada, demanda aparente e corrente de projeto é necessária para evitar equivalências indevidas entre grandezas. Potências expressas em quilowatts (kW) e quilovolt-ampères (kVA) não correspondem à mesma grandeza, assim como a corrente de projeto, em ampères (A), depende da potência aparente e do nível de tensão. Para uma alimentação trifásica equilibrada, a relação entre essas grandezas pode ser expressa por
\begin{equation}
 I = \frac{S}{\sqrt{3}\,V_{LL}},
 \label{eq:corrente-trifasica}
\end{equation}
na qual $I$ representa a corrente de linha, em ampères; $S$, a potência aparente trifásica, em volt-ampères; e $V_{LL}$, a tensão entre fases, em volts. Essa relação não substitui os procedimentos de dimensionamento de condutores e dispositivos de proteção, que também dependem das condições de instalação, da capacidade de condução de corrente, da queda de tensão e das correntes de falta previstas.

\section{Componentes da entrada coletiva e critérios técnicos}
\label{sec:componentes}

A entrada coletiva compreende instalações destinadas à conexão com a rede da distribuidora e à alimentação dos circuitos correspondentes às unidades consumidoras. Sua configuração pode incluir ramal de conexão, ramal de entrada, centro de proteção geral, barramentos, dispositivos de seccionamento, quadros de distribuição e centros de medição. A terminologia e os limites de responsabilidade devem ser observados conforme o documento técnico aplicável, pois a localização física de determinado componente não estabelece, isoladamente, sua propriedade ou a responsabilidade por sua instalação.

O centro de medição concentra os equipamentos utilizados na medição individualizada das unidades consumidoras e pode incorporar módulos de proteção, distribuição e passagem. Sua disposição física influencia as condições de acesso, manutenção, identificação de circuitos e leitura pela distribuidora. A especificação CNC-NDBR-DBR-25-1580-EDBR dedica seções próprias à medição e aos centros de medição, acompanhadas de tabelas e desenhos de arranjos construtivos, conforme \citeonline[seções~7.17 a~7.21]{enel1580}.

A seleção dos dispositivos de proteção e a definição do sistema de aterramento integram o dimensionamento da entrada. A proteção geral deve ser compatível com a corrente de projeto, com a capacidade de interrupção requerida e com as condições de falta. As medidas de aterramento e equipotencialização dependem, por sua vez, das características da instalação e do esquema de alimentação. A especificação CNC-NDBR-DBR-25-1580-EDBR reúne prescrições sobre esses temas nas seções 7.24 e 7.25, sem afastar a necessidade de observância das normas brasileiras correlatas, conforme \citeonline{enel1580}.

\section{Especificações técnicas da Enel utilizadas na pesquisa}
\label{sec:especificacoes}

\subsection{Especificação CNC-OMBR-MAT-18-0126-EDCE}
\label{subsec:0126}

A especificação CNC-OMBR-MAT-18-0126-EDCE, intitulada \emph{Fornecimento de energia elétrica a prédios de múltiplas unidades consumidoras}, foi utilizada pela Enel Distribuição Ceará para estabelecer critérios de projeto e execução das entradas de energia em PMUC. A versão 03, de 31 de março de 2022, é adotada como referência documental anterior nesta pesquisa. O documento abrange condições de fornecimento, medição, proteção, aterramento, cálculo de demanda e documentação necessária à análise dos projetos, conforme \citeonline{enel0126}.

A especificação inclui disposições aplicáveis a situações particulares e critérios condicionados ao local de atendimento. Sua interpretação requer a leitura conjunta do texto, dos desenhos técnicos vinculados ao documento e dos comunicados técnicos cuja aplicação possa ser demonstrada. A existência de disposições regulatórias anteriores a 2025 impede que requisitos já incorporados à versão de 2022 sejam classificados automaticamente como inovações introduzidas pela especificação nacional.

\subsection{Especificação CNC-NDBR-DBR-25-1580-EDBR}
\label{subsec:1580}

A especificação CNC-NDBR-DBR-25-1580-EDBR, intitulada \emph{Fornecimento de energia elétrica em tensão secundária de distribuição até 1 kV -- conexão coletiva}, estabelece requisitos técnicos para o atendimento coletivo no âmbito das distribuidoras da Enel. A versão 2, de 8 de setembro de 2025, registra a substituição da especificação CNC-OMBR-MAT-18-0126-EDCE, entre outros documentos regionais sobre o mesmo tema, conforme \citeonline[item~2]{enel1580}.

Apesar de sua abrangência nacional, a especificação contém condições particulares para as áreas de concessão do Ceará, do Rio de Janeiro e de São Paulo. Sua aplicação exige, portanto, identificar as prescrições compatíveis com o sistema de distribuição e com as soluções admitidas na região estudada. A reunião de requisitos regionais em um único documento não implica uniformidade das tensões nominais, modalidades de conexão ou alternativas construtivas em todas as áreas de concessão.

\section{Vigência e transição das especificações técnicas}
\label{sec:transicao}

A vigência de normas técnicas da distribuidora não se confunde necessariamente com a data registrada em sua folha de identificação. O art. 20 da Resolução Normativa ANEEL nº 1.000/2021 estabelece procedimentos para comunicação e entrada em vigor de novas normas ou de suas alterações. Como regra geral, o §~2º prevê prazo mínimo de 120 dias contado da comunicação aos interessados, ressalvadas as hipóteses excepcionais nele previstas. O §~4º também contempla a preservação de situações constituídas antes da vigência do novo padrão, nas condições regulamentares, conforme \citeonline[art.~20]{aneel1000}.

A data de aprovação de um projeto não identifica, por si só, todas as prescrições aplicadas durante sua análise. Para os fins desta pesquisa, a especificação CNC-OMBR-MAT-18-0126-EDCE constitui a referência normativa anterior, e a CNC-NDBR-DBR-25-1580-EDBR constitui o documento de comparação posterior. A aplicação prática terá como referência a documentação de um empreendimento residencial de múltiplas unidades consumidoras, sem o propósito de reconstituir o histórico administrativo de aprovação ou verificar a execução das instalações.
